\documentclass[12pt,a4paper]{article}
\usepackage{graphicx} % Required for inserting images
\usepackage{chemfig}
\usepackage{lewis}
\author{Svein-Tore Narvestad}
\date{September 2025}

\begin{document}
\textbf{\large{Prøve om kjemiske bindinger:}}\\\\\\
1.	Hva menes med begrepet elektronegativitet og hvor finner vi grunnstoffene som har høy/lav elektronegativitet?\\\\
.\\\\
2.	Hva menes med ioniseringsenergi? Hvilke grunnstoff har høy ioniseringsenergi?\\\\

3. Forklar hvorfor vann er en dipol.\\\\

4.	Hva menes med hydrogenbinding?  Nevn 3 stoff som har hydrogenbinding.\\\\

5.	 Gitt følgende stoffer: $H_2O, NH_3, CO_2,CH_4$ og NaCl.
Hva er elektronegativitetsverdien for grunnstoffene i stoffene over?\\\\

6.	Hvilken bindingstype har vi mellom molekylene i stoffene $H_2O, NH_3, CO_2,CH_4$ og NaCl?\\\\

7.	Tegn elektronprikkstrukturen for stoffene $H_2O, NH_3, CO_2, CH_4$ og NaCl.\\\\ 
8.	Er noen av molekylene dipoler? Begrunn svaret.\\\\
  
\end{document} 
